\chapter{Branchless Execution Kernels}
\section{The Jitter Problem}
Data-dependent branching is the enemy of determinism. In conventional token-based replay, checking enablement involves conditional logic:
\begin{lstlisting}[language=C]
if ((marking & in_mask) == in_mask) {
    marking = (marking & ~in_mask) | out_mask;
}
\end{lstlisting}
If the log is highly variable, the CPU branch predictor fails frequently, causing execution pipelines to flush and introducing nanosecond-scale jitter.

\section{Branchless Mask Calculus}
DTEAM replaces all conditional execution with bitwise selection and boolean masking. Using Bit-Compressed Integer Representation (BCINR) principles, the transition firing kernel is reduced to unconditional arithmetic:
\begin{equation}
M' = (M \ \& \ \neg I) \ | \ O
\end{equation}
Where $M$ is the current marking, $I$ is the input arc mask, and $O$ is the output arc mask. 

To determine enablement and quantify missing tokens without a single branch, DTEAM employs hardware-accelerated population counts:
\begin{equation}
\text{Missing} = \text{popcount}(I \ \& \ \neg M)
\end{equation}
This operation maps directly to the \texttt{POPCNT} CPU instruction. The result is a replay engine that executes in identical time regardless of whether the transition is fully enabled, partially enabled, or completely disabled. 
